\section{Tóm tắt nghiên cứu}
\label{sec:abstract}

Nghiên cứu này thuộc lĩnh vực \textit{Machine Learning} ứng dụng trong Hệ thống Năng lượng Thông minh (\textit{Smart Energy Systems}), tập trung giải quyết bài toán Dự báo Phụ tải Điện Ngắn hạn (\textit{Short-Term Load Forecasting} -- STLF) theo từng giờ trên lưới điện khu vực của \textit{Pacific Gas and Electric Company} (PG\&E) tại California. Dự báo phụ tải chính xác đóng vai trò sống còn trong việc lập kế hoạch vận hành kinh tế, điều độ công suất nguồn, tối ưu hóa phát điện dự phòng (\textit{spinning reserve}) và đảm bảo an ninh lưới điện trước biến động thời tiết cực đoan.

Gần đây, tại cuộc thi \textit{IISE PG\&E Energy Analytics Challenge 2025}, công trình của Roy và cộng sự~\cite{roy2025hourly} (arXiv:2505.11390) đã chứng minh rằng các mô hình học máy dạng bảng phân đoạn theo giờ (24 mô hình XGBoost độc lập cho 24 giờ trong ngày) có khả năng vượt trội hơn các kiến trúc \textit{Deep Learning} phức tạp như \textit{Transformers} (PatchTST, Informer, Autoformer). Mặc dù đạt kết quả xuất sắc, phương pháp của bài báo gốc vẫn tồn tại ba hạn chế lớn:
\begin{enumerate}
    \item Phụ thuộc hoàn toàn vào các biến trễ thời gian thống kê của các thành phần \textit{Principal Component Analysis} (PCA) mà thiếu tri thức vật lý công trình (\textit{thermodynamics inductive bias}), chưa phản ánh được ngưỡng phi tuyến tính trong hành vi kích hoạt hệ thống làm mát/sưởi ấm (\textit{HVAC});
    \item Sử dụng các biến phân loại nhị phân rời rạc theo tháng (\textit{one-hot calendar dummies}) dẫn đến sự mất liên tục tại ranh giới chuyển giao giữa các chu kỳ mùa;
    \item Áp dụng một họ mô hình cây đơn lẻ cho mọi khung giờ mà chưa khai thác tính đa dạng sai số (\textit{error diversity}) giữa các lớp kiến trúc khác nhau dưới các chế độ vận hành lưới điện (\textit{grid operational regimes}).
\end{enumerate}

Để giải quyết triệt để các khoảng trống trên, nghiên cứu này đề xuất một khung phương pháp mở rộng mang tính hệ thống kết hợp ba trụ cột:
\begin{enumerate}
    \item \textbf{Physics-Informed Feature Engineering:} Tích hợp các biến nhiệt động lực học \textit{Heating Degree Days} (HDD) và \textit{Cooling Degree Days} (CDD) dựa trên ngưỡng kích hoạt $18^\circ\text{C}$ chuẩn \textit{ASHRAE}~\cite{ashrae2021fundamentals}, kết hợp cùng chuỗi điều hòa \textit{Fourier harmonics} ($\sin, \cos$) liên tục nhằm mô tả trơn tru chu kỳ nhật triều và chu kỳ mùa thiên văn của Trái Đất.
    \item \textbf{Dual-Regime Discovery:} Khám phá và lượng hóa tính bù trừ sai số giữa mô hình tuyến tính ElasticNet (với L1/L2 regularization)~\cite{zou2005regularization} và mô hình phi tuyến XGBoost~\cite{chen2016xgboost}. Phân tích tương quan phần dư (\textit{residual correlation}) cho thấy mức tương quan thấp ($r = 0.6131$), khẳng định tính đa dạng sai số (\textit{error diversity}) rất cao. Cụ thể, ElasticNet vượt trội hơn XGBoost trong chế độ phụ tải nền ban đêm (\textit{nocturnal baseload}, 00:00--06:00: RMSE $131.33$~MW so với $141.37$~MW), trong khi XGBoost áp đảo hoàn toàn trong các đợt nắng nóng cực đoan (\textit{summer heatwaves}, CDD $> 3^\circ\text{C}$: RMSE $244.16$~MW so với $313.27$~MW).
    \item \textbf{Context-Aware Stacking (CAS):} Đề xuất kiến trúc \textit{Stacking Ensemble} thế hệ mới, trong đó mô hình \textit{meta-learner} (LightGBM~\cite{ke2017lightgbm}) không chỉ nhận đầu ra dự báo của các mô hình cơ sở (\textit{base models}) mà còn được điều hòa đồng thời bởi vector ngữ cảnh khí tượng và thời gian thực (\textit{real-time environmental context}). Cơ chế này cho phép hệ thống tự động định tuyến (\textit{dynamic routing}) và phân bổ trọng số tin cậy tối ưu giữa các mô hình theo từng trạng thái môi trường.
\end{enumerate}

Thí nghiệm được thực hiện trên tập dữ liệu thực tế kéo dài 3 năm (2020--2022) của PG\&E với 5 trạm quan trắc khí tượng (Zenodo: 10.5281/zenodo.17085273~\cite{pge2025zenodo}). Trên tập kiểm tra độc lập (Year 3 / 2022, 8,760 giờ), mô hình CAS đề xuất đạt kết quả vượt bậc với RMSE \textbf{141.81~MW}, MAPE \textbf{4.75\%}, và $R^2$ \textbf{0.9136}. So với baseline tái lập từ bài báo gốc (RMSE $170.84$~MW, MAPE $5.38\%$), mô hình đề xuất giúp cắt giảm sai số tới \textbf{$-$29.03~MW} (tương đương giảm \textbf{$-$17.0\% RMSE}) với ý nghĩa thống kê cao qua kiểm định Diebold-Mariano ($p < 10^{-5}$). Phân tích bóc tách thành phần (\textit{ablation study}) và giải thích mô hình qua biểu đồ nhiệt \textit{SHAP} 24 giờ~\cite{lundberg2017shap} khẳng định vai trò then chốt của các đặc trưng vật lý và cơ chế định tuyến ngữ cảnh trong việc nâng cao độ chính xác và độ tin cậy của hệ thống lưới điện thông minh.
